\documentclass[11pt]{article}
\usepackage{amsmath,amssymb}
\usepackage{iftex}
\ifPDFTeX
  \usepackage[T1]{fontenc}
\else
  \usepackage{unicode-math}
  \setmainfont{STIX Two Text}
  \setmathfont{STIX Two Math}
\fi
\usepackage[margin=1in]{geometry}
\usepackage{microtype}
\emergencystretch=3em
\usepackage{xurl}
\usepackage[authoryear,round]{natbib}
\usepackage{hyperref}
\hypersetup{hidelinks,pdftitle={A Computer without a First Step?}}

\title{A Computer without a First Step?}
\author{}
\date{Draft for review. Not submitted.}

\begin{document}
\maketitle

\begin{abstract}
An invented dossier of three laboratories has one consistent record for every local bit-function, and still no laboratory is first. That arrangement is easy to mistake for a theory of computers. Universal consistency makes the local substitutions admissible together. It does not supply an efficient way to find the record, a physical implementation, or warrant that such a machine is metaphysically possible. Finite explicit tables can be searched by brute force, so the gap is not uncomputability. A reversible extension, tested by every local unitary on matched finite ancillas, draws the same line. Ordinary circuits that loop through time are a different arrangement. The proposal is a two-level assessment of a candidate computer, not a definition of computation.
\end{abstract}

\noindent\textbf{Keywords.} computational implementation; consistency; causal order

\section{An answer without a first message}
\label{sec:open}

The following scene is invented. It is not a report of a laboratory, a patent, or a machine anyone is asked to build.

An archivist is handed a dossier on three workshops. Each workshop has one incoming bit and one outgoing bit. Each technician may choose, independently, any of the four functions from a bit to a bit: copy the bit, flip it, ignore it and write \(0\), or ignore it and write \(1\). There are sixty-four such choices. The dossier does not say which technician speaks first. It lists, for each of the sixty-four choices, a single triple of outgoing bits that fits the global rule. An engineer, who has spent a career laying cables in an order, asks the natural question: which workshop must send the first message? The archivist does not draw a route. The archivist points at the list.

The engineer is not being fussy. Ordinary computers do have a first step, or at least a first layer, once the diagram is unrolled in time. A gate takes its inputs from something earlier. A program counter moves. If one has only ever seen machines of that kind, it is reasonable to think that determinate local operations require a first laboratory in one fixed order. Call that the first-step principle. The principle is not a quotation from any of the accounts discussed later. It is the inference the scene is built to test. Those accounts, as will be plain, are stronger and more careful than the principle.

The thesis of this paper is narrow. Universal consistency, in the sense of one admissible record for every local substitution of the relevant kind, is a condition on which operations may be composed. It is not an algorithm for producing the record, not a physical implementation, and not a bridge from the formalism to metaphysical possibility. A reversible quantum extension of the same tables, under a stated family of tests, does not restore a first laboratory and does not supply the missing implementation. A metaphysics of computers should therefore keep two questions apart. One question is whether the local pieces are admissible together. The other is whether anything physically organized computes. Neither question answers the other.

``Without a first step'' means only that the dependence among the three workshops need not pick out a laboratory that is first in a fixed order. It does not mean that nothing happens in succession inside a workshop. A technician can still do one thing and then another. The clocks on the walls are not the subject. The subject is whether the global rule has to begin at one of the three doors.

The engineer’s question will be left standing. Sections~\ref{sec:leibniz} and~\ref{sec:turing} say why a consistent description, and even a computable finite description, is not yet a machine. Section~\ref{sec:dossier} returns to the dossier and shows what the sixty-four records do and do not show. Section~\ref{sec:ledger} asks what a reversible ledger adds. Section~\ref{sec:seller} separates a unique answer from a way of finding it. Section~\ref{sec:impl} puts the two-level proposal against causal and mechanistic accounts of implementation. Section~\ref{sec:obj} answers the objections that the dossier is a mere fiction, that no serious theory ever required a first laboratory, and that a quantum ledger quietly restores a classical computer underneath.

One boundary is reserved for a note, because it is easy to inflate.\footnote{With infinitely many parties, pairwise exclusivity of histories still yields at most one fixed point, while existence can fail. The general question whether every exclusive map of Baire class one has a fixed point is open. Nothing in the argument uses that question. The parties and alphabets below are finite.} The mathematical claims are reported in an unpublished technical note. Related drafts of that note are already public. This manuscript does not claim that anonymity is secure, and it gives no address.

\section{Inspecting the parts}
\label{sec:leibniz}

Leibniz’s \emph{Monadology} is not a treatise on computers, and nothing below treats it as one. The text used here is Jonathan Bennett’s modernized English rendering of the 1714 essay. Bennett marks his own additions and updates the vocabulary. Section numbers refer to that rendering. No sentence is quoted. The use of the essay is a comparison of explanatory demands, not a claim of historical influence on any later formalism.

In \S 17 Leibniz asks the reader to imagine a machine that thinks, perceives, and has experiences, enlarged until one could walk inside it as through a mill. The tour would show parts pushing parts. It would not, he holds, show perception. The scene is a thought experiment. It is not an incident at court, and it is not about causal order among laboratories. Its target is the unity of perception. The only moral taken here is preparatory and negative. A walk through the parts is one task. An account of what the organized whole is doing is another. The paper does not conclude anything about consciousness, and it does not use the mill as evidence that mechanism fails in general. Piccinini’s later account, discussed in Section~\ref{sec:impl}, is a mechanistic account of computing that is untouched by a remark about perception.

Sections 31--38 draw a different line, and it is the line this paper needs. Leibniz distinguishes truths whose opposite is contradictory from truths whose opposite is possible. The former, on his telling, can be analyzed until one reaches identities. The latter require a sufficient reason, and the train of reasons runs into ever more minute detail, without a last step inside the series of contingencies. He does not say that a consistent description is already the reason the thing exists. Consistency keeps a description from destroying itself. A reason, in his sense, would say why this consistent description, rather than another, is the one that obtains. That is a demand, not a theorem of later physics. It is also not an identification of consistency with actuality. The opposite reading would erase the distinction he is drawing.

Sections 53--59 press the demand onto the choice among possible universes. There are, he holds, many possible universes in God’s understanding, and there must be a reason for the selection of one. His reason is theological: a principle of fitness, aiming at the greatest variety together with the greatest order. The theology is not adopted here. What is adopted is the shape of the question. After a description has been shown to fit, one may still ask why that description is realized. A table of records can answer the first question and leave the second untouched. The refusal, later in the paper, to read a metaphysical computer off a finite table is a refusal of Leibniz’s own inference from the demand to a divine selection. The demand can be kept while the inference is dropped.

Sections 78--81 coordinate soul and body by harmony rather than by influence. The two realms, of final causes and of efficient causes, follow their own laws and agree. The comparison with a process function has to be drawn against the grain. A process function is not a population of windowless substances. It contains operational dependencies: what one party receives is a function of what others emit. And it is defined by a license to substitute local functions. Harmony, in Leibniz’s picture, is coordination without that sort of mutual operational dependence. The dossier in Section~\ref{sec:open} is full of dependence. The resemblance is only this: both pictures refuse to explain a global fit by a single push from one part to the next in a line. The difference is that Leibniz’s fit is secured by a prior decree, whereas a process function is a constraint on which global rules are admissible given free substitution of local maps. One should not say that the dossier proves pre-established harmony, or that monads were an early theory of indefinite causal order.

The engineer’s question survives this history. Leibniz asks for a reason that a consistent whole obtains. The archivist has so far offered only the whole. Walking through three workshops, like walking through the mill, would show bits and functions. It would not by itself show why those functions have a unique joint record, and it would not show that any cable was ever laid.

\section{The person at the tape}
\label{sec:turing}

Turing’s 1936--37 paper, with its 1938 correction, is about numbers and mechanical procedures, not about laboratories wired into a cycle. The relevant contrasts can be stated without treating the paper as a description of a laptop.

A computing machine, in the opening sections, has finitely many configurations. At each step it scans one square of a tape. A transition rule says how the configuration and the scanned symbol determine a written symbol, a move of the tape, and a next configuration. The restrictions are motivated by the limits of a human computer: a bound on how much can be discriminated and attended to at once. That is a theory of effective steps. It is not a theory of a device whose global state is solved all at once.

Turing distinguishes numbers that can be defined from numbers that can be computed. A law may pick out a real number uniquely and still fail to supply a mechanical procedure that writes its digits. The universal machine does not close that gap by magic. It is given a description of another machine and simulates that machine’s steps. Simulation is a formal result about descriptions. It does not say that every mathematical structure is a piece of hardware, and it does not say that a uniquely specified table is already being computed by the world in which the table is written down.

One piece of vocabulary should be fenced off. Turing calls some machines circular and others circle-free. The distinction concerns whether a machine prints infinitely many output figures or eventually stops producing them. It has nothing to do with a physical loop in which an effect is among the causes of its own past. Using the word ``circular'' as a bridge from his paper to causal loops would be a pun.

The qualification that matters for the dossier is the opposite of an undecidability claim. The three workshops have finite alphabets, and the global rule, once written down explicitly, is a finite table. For any chosen triple of local functions, a programmer can list the eight possible output triples and test which of them, fed through the rule and the local functions, returns itself. The search ends. Turing’s theorem that there is no general mechanical decision procedure for all machines does not say that this finite search fails. What the search does not provide is an account of cost that survives changes in how the table is encoded, how large the alphabets grow, or how the answer is delivered to someone who cannot see the whole dossier. Those are questions about resources and access. They are not questions about whether a finite explicit instance is computable.

A further distinction keeps ordinary computers from being drafted into the wrong camp. A thermostat, or a recurrent network, has a loop in its diagram. The air warms, the sensor trips, the furnace stops, the air cools, the sensor clears. Unrolled, the history is a sequence of states of the same devices through time. There is a first moment of the episode even if the diagram draws an arrow in a circle. The succession is inside the life of the device. The dossier is not like that. Its three parties are not three times. They are three slots, and the rule may make each slot’s input depend on another slot’s output, with no slot earliest in the dependence. A metaphysics that treats every circular arrow as the same phenomenon will confuse a repeated event with an all-at-once composition of parties. The first-step principle is plausible for the thermostat because the thermostat’s loop is already a temporal succession. It is not thereby plausible for the dossier.

Turing’s human computer works in succession. The archivist’s rule is a compatibility condition on completed records. Both can be exact. They are not the same job. A serial simulation can discover which record fits without the three workshops ever having been wired together. Conversely, a rule that admits exactly one record for each substitution does not, by itself, describe a sequence of steps from an arbitrary seed. That second point gets its own scene in Section~\ref{sec:seller}. The first point is already enough to stop a slide from ``the record is determined'' to ``a computer of the ordinary sort is present.''

\section{Sixty-four records}
\label{sec:dossier}

Return to the invented dossier, and give the archivist’s rule a definite form. It is the three-bit rule often called the Lugano rule in the literature on classical processes without a fixed causal order \citep{baumeler2016,baumeler2019}. Write the outgoing bits as \((a,b,c)\) and the incoming bits as
\[
(\lnot b \land c,\; \lnot c \land a,\; \lnot a \land b).
\]
Each incoming bit depends on the other two outgoing bits. There is no party whose incoming bit is independent of the others, and there is no ordering of the three parties in which each depends only on parties earlier in the list.

A local intervention is a triple of functions, one per party, from that party’s incoming bit to that party’s outgoing bit. A record is a triple of outgoing bits that, sent through the rule and then through the chosen functions, yields itself. The rule is a process function when every intervention has exactly one record \citep{baumeler2016}. The unpublished note contains the general finite-alphabet definition and the checks. One intervention can be seen by hand. Let every party copy its input. A record must then satisfy \(a = \lnot b \land c\), \(b = \lnot c \land a\), and \(c = \lnot a \land b\). If \(a = 1\), then \(b = 0\) and \(c = 1\), and the third equation says \(c = 0\). If \(a = 0\), then \(b = 0\) and \(c = 0\), and \((0,0,0)\) satisfies all three. Exactly one record. A second intervention shows that the record is not a slogan repeated under every choice. Let every party ignore its input and write \(1\). The only triple the local functions can emit is \((1,1,1)\). The rule sends \((1,1,1)\) to \((0,0,0)\), and the constant functions send \((0,0,0)\) back to \((1,1,1)\). So \((1,1,1)\) is a record, and nothing else can be, because nothing else is in the image of the local functions. The identity choice and the constant choice have different records. The archivist is not pointing at one favored history and calling it inevitable. The other interventions are a finite list of the same kind. The two cases on the page do not prove the universal claim. They show what the claim says. The universal claim is that the list has one success for every entry. That claim is part of the standing example, not a new enumeration offered here.

What the scene tests is the first-step principle. The local operations are as determinate as bit-functions can be. The joint record is unique for each choice. No laboratory is first in the dependence. So determinate local operations, plus a unique global record for every substitution, do not require a first laboratory in a fixed order. The principle, stated as a general rule about any such arrangement, is false inside this formalism.

What the scene does not test is the engineer’s capacity to install the wires. Nothing in the sixty-four lines is a spacetime diagram, a parts list, or a procedure for setting the technicians’ functions by hand in a room. \citet{baumeler2019} show that deterministic classical dynamics compatible with arbitrary local operations, including arrangements in which three parties are each in the past and the future of the others, admit a reversible realization in their model. ``Possible,'' in that result, is possibility relative to the classical process theory. It is not a report that the wires have been laid. The archivist’s answer is complete as an answer about records. It is silent as an answer about construction. The engineer’s question remains the question of the paper.

Two failures of the fixed-point count should not be collapsed. If an intervention has no record, the local choice and the global rule cannot be satisfied together. That is the flavor of contradiction: the demands exclude every candidate. If an intervention has two records, both candidates fit. The local choice has not picked which history occurs. That is underdetermination, not a contradiction. A law can leave a parameter free without being inconsistent. The process-function condition excludes both failures, by a policy of determination: exactly one record, for every substitution. The policy is stronger than logical consistency alone. Someone who says ``consistency is not determinism'' is right about the words. The dossier uses a special normalization. It should not be described as the discovery that multiple solutions are metaphysically impossible. They are excluded from this class of rules because the class was defined to exclude them.

The same policy is what makes the quantifier do philosophical work. A story that checks one intention, and finds one fitting history, has not faced the substitutions. The sixty-four lines are the insistence that every locally available function is on the list, including the boring constant functions and the bit flip. Admissibility is universal or it is not the condition under discussion. A rule that works for the functions a narrator finds natural, and breaks for a function nobody bothers to try, is not a process function. The archivist’s smugness, if any, is licensed only by the whole list.

\section{What the ledger buys}
\label{sec:ledger}

A second invented scene, continuous with the first. The archivist now replaces the bit-functions with coherent quantum operations. Each party may insert a unitary on the incoming system together with a private finite ancilla, provided the dimensions match so that a unitary of that shape exists. The archivist keeps a ledger: a source system and a sink system, and a reversible map that writes the global rule into the incoming registers and stores the outgoing string in the sink. The engineer asks, again, how the workshop is built. The archivist answers by displaying the ledger.

The display is not yet the quantum condition. For an arbitrary rule \(w\), not only for a process function, there is a permutation of a finite basis,
\[
U\lvert o\rangle\lvert e\rangle = \lvert w(o)+e\rangle\lvert o\rangle,
\]
with addition in each party’s input alphabet. A permutation of a finite basis is unitary. The ledger exists for rules that have no record at all, and for rules that have many. Storing outputs reversibly is bookkeeping. It does not say that the local substitutions compose into a unitary on the source and the ancillas.

The condition that does say that is narrower, and the narrowing is the result. Take the canonical extension \(U\) just written. Insert arbitrary local unitaries and arbitrary finite ancillas of matching dimension. The closed loop is unitary for every such insertion if and only if \(w\) is a process function. One direction uses a combinatorial fact. If two outgoing strings differ and \(w\) is a process function, some party has a different output symbol and the same input symbol on those two strings. The Gram matrix of the closed loop factorizes over parties, and that party’s factor vanishes, so the rows that should be orthogonal are orthogonal. The dimensions make the loop square, hence unitary. The other direction is a counter. If some local functions have \(N\) records and \(N\) is not one, those functions can be written as a permutation on a larger space, and a column of the loop then has squared norm \(N\). An operator with the wrong column norm is not unitary.

The equivalence is an equivalence for this canonical extension and this family of tests. It is not the claim that every unitary dilation of every kind is a process function, and it is not the claim that purification in any wider sense coincides with consistency. A reader who compresses it into ``quantum reversibility if and only if consistency'' has dropped the restriction that makes the sentence true. Column norms and row orthogonality are different tests. One failing insertion rejects a rule. The forward direction requires every insertion to pass.

Reversibility, in this construction, is not the demand that \(w\) itself be a one-to-one function of the outgoing string. The Lugano rule is not. The sink holds the outgoing string, and the source, read backwards, holds the difference between the incoming string and \(w(o)\). The inverse uses the ledger. It does not force the classical rule to be a permutation on its own alphabet. That is a correction to a picture in which only bijections are reversible. It is still a correction inside the formalism. The source and the sink are extra systems in the mathematics. They are not a discovered power supply.

What the ledger buys, then, is the removal of one formal obstruction. If the classical rule already has one record for every local function, this canonical extension does not become non-unitary when the parties insert unitaries. If some local function has the wrong number of records, a unitary insertion exhibits the failure. Quantum theory, in this limited sense, is not an extra filter laid on top of the fixed-point condition. The filter is the condition. The engineer’s question is not the filter. Nothing in the equivalence is a Hamiltonian, a preparation and readout protocol, a noise model, or a region of spacetime. \citet{vilasini2024} give no-go results about embedding cyclic information-theoretic structures in acyclic spacetimes, under assumptions that are not re-argued here. The ledger is not an exception to those results, because the ledger is not an embedding.

The equivalence is not a recoding of the same sentence in a longer alphabet. The permutation \(U\) is available before anyone asks whether \(w\) is a process function. A local unitary that is not diagonal in the computational basis mixes basis states the classical table treated as distinct alternatives. If, after that mixing, the map from source and ancillas to sink and ancillas fails to be unitary, the fixed-point clause has not survived the change of test. The theorem says that for a process function it does survive, and that when it survives for every such test the clause was in force. Survival is a result, not a definition. The ancillas are part of the stated test, not a decorative extra. The dimension count that lets a local map be unitary uses them, and the converse gadget uses an ancilla to count records. Dropping the ancillas would be a different theorem. Stochastic classical maps are a different test again: a valid probabilistic process need only assign total probability one under every product of local stochastic maps. Unitary insertions, instruments, and deterministic bit-functions are three quantifiers. The equivalence above is the middle one, for the canonical lift. The classes are worth listing, because a shift among them is a change of subject. Deterministic bit-functions are the dossier’s substitutions: four maps on one bit, sixty-four triples, one record each. Stochastic maps are mixtures of those functions, and a valid probabilistic process is one that still assigns total probability one when the parties randomize. Local unitaries on matched finite ancillas are the ledger’s substitutions, and they are not a notation for the bit-functions. Instruments that are not unitary are a further class: a local completely positive map can be dilated to a unitary on a larger finite space and the environment then discarded, which is how the note passes from the unitary tests to ordinary instruments. The equivalence in this section is only the third class, for the canonical \(U\). A result about bit-functions is not a result about instruments, and a result about one favorable unitary is not the universal test. The archivist who answers the engineer’s question about cables by quoting the bit-function count, and then, when pressed, quotes the unitary equivalence as if it were the same sentence, has changed the intervention class. The paper’s discipline is to name the class every time the quantifier is used.

Saying so is what keeps the archivist from changing the subject halfway through the answer.

\section{The promised answer}
\label{sec:seller}

A third invented scene. A seller, shown the dossier, offers a service. Because each of the sixty-four choices has exactly one completed record, the seller’s machine will print that record at once when a choice is named. No search, the seller says. The uniqueness is the printing. A programmer across the table writes a different service. The programmer lists candidate output triples and tests the fixed-point equation. For eight candidates and a fully written table, the program stops quickly. The programmer does not claim to have instantiated three cyclic workshops. The programmer claims to have evaluated a finite predicate.

The seller’s inference is the one to reject. A unique solution is not an efficient dynamics, and it is not a dynamics at all. The programmer’s method shows the complementary point. For a finite explicit table, brute force is a method. The absence of a closed-form shortcut is not the absence of computability. Anyone who reaches for Turing’s undecidability theorem at this moment has changed the subject from a finite table to a general decision problem. The interesting remainder is not ``can the record be found?'' It is ``what resources, what encoding, and what access does the finding use, and do those coincide with the organization of the workshops?''

A smaller diagram, labeled as not a process function, separates uniqueness from iteration. Let \(F(0) = 0\), \(F(1) = 2\), and \(F(2) = 1\). Then \(0\) is the only fixed point. Start at \(1\) and iterate, and the sequence is \(1, 2, 1, 2, \ldots\). It never arrives at the unique solution. Uniqueness of a solution does not mean that repeated application from an arbitrary seed converges to it. \(F\) is not a process function. A process function has to survive every local substitution, which is a stronger demand than having one fixed point as a self-map of a three-element set. The diagram is only an illustration of the illicit slide from ``exactly one solution'' to ``relaxation will find it.'' The dossier is vulnerable to the same slide, and the diagram shows the slide in a place where the arithmetic fits on a line.

The seller may retreat to a different literature. There is a model, following Deutsch, in which a closed timelike curve is treated as a fixed point of a channel, and in which the computational power of that model can be analyzed. \citet{aaronson2008} show, from the assumption that such curves exist and from their stated model of them, that quantum and classical computers in the model would both have the power of polynomial space: the problems solvable by an ordinary computer with a polynomial amount of memory. That is a large power, and it is not hypercomputation. It is also not a theorem about process functions. Their fixed points are fixed points of channels, with a correctness condition that ranges over consistent fixed points of a different shape, in a uniform circuit model. The dossier’s condition is one deterministic record for every local function on finite alphabets. Importing polynomial-space power from one theory to the other would require a translation nobody has given, and this paper does not offer one. The citation is a fence. A fixed-point theory of computation can have resource consequences. Those consequences stay inside the theory that proved them.

A related fence concerns simulation by ordinary circuits. Some families of quantum processes with indefinite causal structure can be simulated by ordered circuits under stated access assumptions. The existence of a simulation in a given model is a result about that model’s gates, oracles, and allowed dilations. It is not a corollary of the sentence ``the table has one record.'' Nor does the failure of a particular simulation follow from cyclic dependence alone. The question ``does this compute faster?'' is a question about an access model. The archivist has not answered it. The programmer, who merely searches eight lines, has answered a different question, and has answered it without building the workshops.

Access is the part the seller skips. The programmer’s search begins only when the table is in hand: the eight lines, the rule, and the chosen functions. If the rule is what one is trying to discover, the search has not started. If the alphabets grow, the list of candidates grows with them, and “finite” does not say whether the list fits in the afternoon. If the output must be delivered to someone who cannot see the dossier, delivery is a further task. None of these remarks is a complexity classification. They are the reason a counting argument is not an account of computation. Uniformity, encoding, and the difference between holding the table and being the workshops have to be specified before anyone is entitled to a hardness claim, and this paper states none.

The positive implication is modest and practical. Operational access and resource accounting are explanatory jobs of their own. A metaphysics of computers that stops when the records are unique has not yet said who can read the record, how the table is supplied, or what is consumed in producing the answer. The seller’s instant printer is a picture of those jobs being declared finished by a counting argument. The picture is empty. The programmer’s search is a picture of one job being done. It is not a picture of the engineer’s job.

\section{Implementation and modal warrant}
\label{sec:impl}

The three scenes score differently, which is why they were invented as a set rather than as ornaments.

The dossier passes level A for deterministic bit-functions and fails to answer level B. The seller’s printer does not pass level A or level B. It misreads a count as a dynamics. The programmer’s search is a computation in the ordinary sense: a person or a machine, given the table, evaluates a finite predicate. That computation does not implement the three workshops. It implements a search. The ledger passes level A for the canonical unitary tests exactly when the rule is a process function, and it still fails to answer level B. A theory that gives all three scenes the same grade has not been reading them.

Piccinini’s treatment of malfunction marks the same cut from the other side. A computing mechanism can miscompute: it is the mechanism whose function is to produce a certain string, and it produces another. The failure is a failure of a device that remains that device. An inadmissible rule, one with no record or with two, is not a computer having a bad day. It is not in the class. Digestive systems, on his comparison, are mechanisms that are not computers at all. The dossier’s exclusions and a broken printer’s errors are different kinds of “does not compute.” Keeping them apart stops a metaphysics from treating every formal defect as a hardware fault, and every hardware fault as a proof that the formalism was inconsistent.

The two-level proposal can now be stated as a constraint on theories, not as a definition of computation.

Level A is compositional admissibility. The local operations are specified, the substitutions are specified, and the global rule is admitted only when every substitution has the right number of histories. For the dossier, the number is one, and the substitutions are all deterministic bit-functions. For the ledger, the substitutions are all local unitaries on matched finite ancillas, and the canonical extension is admitted only when every such substitution induces a unitary loop. Level A can be satisfied by a cyclic dependence. It can be checked, in the finite explicit case, by a search. It does not say how the check is implemented in hardware, and it does not say that the rule is the function of a physical device.

Level B is an implementation and access story. It says which physical parts correspond to which formal slots, under what manipulations those parts would do the corresponding thing, how the output is read, and what has to be stable for the correspondence to be the device’s behavior rather than a mapping imposed after the fact. Level B is where causal, counterfactual, and mechanistic accounts of concrete computation do their real work. The proposal is that a theory errs if it treats level A as a substitute for level B, and errs differently if it treats a first laboratory in a fixed order as a requirement on level A merely because familiar circuits have a first layer. The proposal is not a necessary and sufficient analysis. It is a discipline for keeping the engineer’s question and the archivist’s answer from being scored as the same exam.

\citet{chalmers1996} gives the causal-combinatorial standard against which this discipline should be measured. He is answering Putnam’s claim that every ordinary open physical system implements every finite-state automaton. The reply has several layers, and they should not be flattened into the first-step principle. First, implementation conditionals are not material conditionals about one actual run. If the system is in a physical state that maps to a formal state, it must transit to a physical state that maps to the successor, in a way that supports counterfactuals: however it got into the first state, within the relevant range of circumstances, it moves to the second. Unexhibited transitions count. A mapping built only from the states that happened to occur has not implemented the automaton. Second, inputless automata are too weak to carry the argument. A clock and a dial, Chalmers argues, can be mapped onto every inputless automaton while satisfying the strong conditionals, because the state space of such an automaton is essentially an unbranching sequence into a cycle. Inputs create branching, and branching is what makes the mapping hard to fake by disjunctions of unrelated states. Third, for systems whose states have internal structure, he introduces combinatorial state automata. An internal state is a vector of independently variable components. He requires those components to correspond to distinct physical regions, and he requires the causal transitions of the physical system to mirror the formal transitions of the vector, including the local dependencies. A finite combinatorial automaton is no more powerful, in the computability sense, than a finite-state automaton. The implementation conditions are stricter, because the parts have to vary independently and the conditionals have to hold componentwise.

Two consequences matter. The process-function quantifier resembles, at the formal level, Chalmers’s demand that every transition be mirrored, not only the transition that appeared in one story. The sixty-four lines are a robustness test under substitution of local functions. That resemblance is real, and it is why level A is not a triviality. The resemblance stops before implementation. Chalmers’s conditionals are conditionals about a physical system: if it were in this physical state, it would go to that one, because of how it is causally put together. The dossier’s substitutions are substitutions of formal functions on bits. Nothing yet shows that the three workshops are independently manipulable components, or that the mapping would survive an intervention on a physical part rather than a rewrite of a line in the file. To equate the two is to treat a slot in a definition as a dial one can turn. Chalmers’s own account is a warning against that equation, not an invitation to it.

The other consequence is that Chalmers does not hold the first-step principle. A combinatorial automaton can have rich local dependencies. Nothing in the definition requires that the components be arranged in one global order with a unique earliest component. Cellular automata and Turing machines are among his examples, and they are full of neighbor-to-neighbor dependence. The false implementations he still worries about, late in the paper, are astronomical systems that satisfy the letter of the conditionals without being the sort of machine anyone would build. He leaves the account unfinished on that point. A fair use of his view is therefore not ``causal implementation has been refuted by a cyclic table.'' A fair use is that even a strong causal and localizational account does not collapse into the first-step principle, and does not become optional once a formal substitution test has been passed. The dossier fails to be an implementation, on his standard, until the physical conditionals are supplied. The dossier succeeds in showing that those conditionals are a further demand, not a spelling-out of the fixed-point count.

\citet{piccinini2007} asks a different question, closer to the engineer’s. A computing mechanism, on his account, is a mechanism whose function is to generate output strings from input strings, and perhaps internal states, according to a general rule that applies to the relevant strings. The account is mechanistic in the sense used in biology and engineering: one explains a capacity by the parts, their functions, and their organization. He separates this from a semantic theory. Computation, in the sense he is reconstructing from computability theory and computer practice, does not require the strings to represent. A digit is a concrete counterpart of a letter from a finite alphabet: a state or a part that the mechanism treats uniformly when it is of one type and differently when it is of another. Mechanistic explanation distinguishes the properties that matter to the capacity from the properties that do not, and it distinguishes success from malfunction. Digestion is his comparison class. It is a mechanism. It is not, on the ordinary scientific use of the term, computing.

The dossier can be described in his vocabulary, and the description does not come out as a computer. There are strings, and there is a rule. There is not, in the scene as given, a mechanism whose function is to produce those strings. The archivist’s list is a constraint on which records fit. A constraint can be exact without being the function of an organ or an artifact. Piccinini’s objectivity condition, that it is a matter of fact which computations occur, is a reason to refuse a mapping that is imposed by the archivist’s choice of description alone. His explanation condition is a reason to ask how the parts and their organization yield the capacity. The engineer is asking that question. The two-level proposal says that level B is where it is answered, and that level A does not become false because the answer is still missing. A cyclic admissible rule is a real formal object. It is not yet a computing mechanism in Piccinini’s sense.

A later development should be marked without being counterfeited. \citet{anderson2024} defend what they call a robust mapping account: a physical system implements a computing system when it satisfies a robust computational description and thereby bears the physical signature of that computation. The chapter that states the account was not available in full for this draft. The published abstract of that chapter is the text relied on, together with Piccinini’s 2007 paper. Nothing below attributes a specific information-theoretic identity to them beyond the abstract’s claim that robustness is the further constraint a mapping needs. The two-level proposal is compatible with that claim. Robustness, usability, and a physical signature are level B conditions. They may absorb the proposal rather than compete with it. The proposal’s only insistence is that a universal fixed-point condition on formal substitutions is not already that signature.

Put the engineer back in the room with Chalmers’s conditional in hand. The conditional does not ask whether a line in a file has a fixed point. It asks whether, were a particular physical part in a particular condition, the device would move to the successor condition, and whether that would remain so for the conditions that did not occur on this occasion. The archivist can rewrite a local function on paper. That rewrite is not yet a manipulation of a part. To turn it into one, someone must say which knob, which register, or which degree of freedom is the incoming bit, what counts as setting it, and what would have happened had it been set the other way. Those are empirical questions about a device. They can be asked of a thermostat, and the answers are boring and physical: the sensor is a part, the set point is a setting, the furnace’s response is the transition. They cannot be asked of the dossier until the dossier is given parts. The first-step principle mistakes the thermostat’s success for a formal law. The archivist’s list mistakes a formal law for the thermostat.

The same gap appears if one tries to save the principle by saying that every computer has loops, so the dossier is just a computer seen from above. The thermostat’s loop is a sequence. At noon the room is cool and the furnace is on. At half past the room is warm and the furnace is off. The later state is not an input to the earlier state. The diagram’s circular arrow is a summary of a succession that, written out, has a first time. The dossier’s dependence is not a summary of three times. It is a rule in which the bit delivered to the first workshop is a function of bits emitted by the other two, and likewise around the table, with no parameter left that orders the three deliveries. One can of course embed the rule in a larger story that does have times: a simulator steps through candidate records on Monday and prints the winner on Tuesday. That embedding has a first step. It is the programmer’s search again, not the three workshops. Confusing the embedding with the rule is how a temporal machine gets credited with having explained an all-at-once constraint.

Level B can be partial without being optional. A laboratory might implement two of the three slots and simulate the third, or implement the classical record while the unitary ledger remains a calculation. Piccinini’s taxonomy of mechanisms by what they in fact do leaves room for such partial stories: a device computes the function it is organized to compute, not every function a theorist can associate with its shadow. The robust-mapping abstract makes the related demand that the physical description be robust, not an after-the-fact alignment. Partial implementation, if the later account’s abstract is a guide, is a deviation from that robustness, not a silent promotion of every associated formalism into the device. The two-level proposal does not need the unpublished details of that deviation to say the following: a simulator of the dossier is a computer of a search or of a circuit that has been shown, in some access model, to reproduce an input-output table; it is not, by that success, the cyclic composition itself.

Modal warrant is a separate mistake, easier to make once the mathematics is clean. A countermodel defeats an entailment inside the axioms that define it. The first-step principle, read as ``any arrangement of determinate local operations with a unique record for every substitution has a first laboratory,'' is false of the Lugano rule. That is a result about the principle. It is not a result about which machines exist in the actual world, and it is not a result about which machines exist in some metaphysical sense broad enough to include every consistent formal object. Leibniz’s demand for a reason why one possible universe is selected is a reminder that consistency was never actuality. Dropping his theology does not install the opposite theology, in which every admissible table is a possible machine. Physical possibility requires an implementation story of the sort level B asks for, and even that story, if it were supplied, would be an empirical and engineering claim. Metaphysical possibility requires whatever further modal premises one adopts. This paper adopts none that would carry a finite table into a possible world. The clean ledger removes an obstruction to unitarity inside a formalism. Removed obstructions are not existence proofs.

What, then, would a metaphysician have to add before saying that the workshops are a possible computer? Not a longer proof of the fixed-point count. The additions sit on level B, and they are ordinary. There must be a proposed decomposition of a physical system into parts that can vary independently in the way the three bits vary. There must be a claim, open to counterexample, about what those parts would do under settings that did not occur. There must be a story of how the output is read without smuggling in a later simulator and calling the simulator the workshops. There must be a decision about malfunction: what the device does when it fails, as opposed to which formal rules were never candidates. None of these is supplied by unitarity of \(U\), and none is forbidden by it. A causal account can take them as its own homework. A mechanistic account can take them as the demand for parts, functions, and organization. The two-level proposal fails if it is read as a rival that declares those homework assignments unnecessary. It succeeds if it stops a formal admissibility result from being graded as though the homework were done.

There is a temptation to hear the cyclic dependence as a grounding cycle: the output is grounded in the input which is grounded in the output. The formalism does not say that. It says that a certain function from outputs to inputs has one fixed point under every local substitution. A dependency graph of a function is not a graph of metaphysical ground. Explanatory direction is not settled by the existence of a unique solution. An equation can constrain a pair of quantities without either quantity being the ground of the other. The dossier constrains records. It does not exhibit a productive cause, and it does not exhibit a sufficient reason in Leibniz’s sense. The explanatory relation actually used in this paper is modest. The fixed-point condition explains why certain combinations of local functions and global rules are inadmissible: they have the wrong number of records. It does not explain why a workshop with that rule exists.

\section{Three objections, and the workshop again}
\label{sec:obj}

A global solution does not explain itself. The fixed-point condition explains an incompatibility: these local functions and this rule cannot be satisfied together, or cannot be satisfied in only one way. That is a constraint on a class of descriptions. A constraint of that kind is not a productive cause of a workshop, not a metaphysical ground of the outgoing bits, and not a sufficient reason why this table rather than another is realized. The seller’s offer treats the count as if it were already the production. Leibniz’s sufficient-reason demand, stripped of the theology, is the reminder that production is a further question. The paper uses the constraint in the first role only.

The first objection says that the dossier is a mathematical fiction, and that fiction cannot discipline a theory of physical computers. The objection is granted in the form in which it is true. The scene is not a physical counterexample. No cable is displayed. If a theorist defines computation as whatever physically realizable systems do, the dossier does not show that any such system lacks a first step. What it shows is that one may not infer a first laboratory from local functionality plus formal consistency alone. The inference needs an extra premise connecting the formalism to a physical organization. That extra premise is level B. Calling the dossier a fiction is a way of agreeing that level B has not been discharged. It is not a way of restoring the first-step principle as a formal truth.

The second objection says that no serious causal or mechanistic account requires an acyclic party graph, so the paper has attacked a straw man. The objection is also granted, and it was built into the exposition. Chalmers’s combinatorial automata are allowed rich dependence. Piccinini’s mechanisms are individuated by function and organization, not by a ban on cycles in a wiring diagram. The first-step principle is a further claim, suggested by the look of ordinary circuit diagrams and by the temporal succession inside recurrent devices. The paper’s demand on the serious accounts is only that they say, in their own terms, what they require beyond level A. If the answer is ``causal conditionals on independently variable parts,'' or ``a mechanism whose function is to produce the strings,'' those answers are level B answers. They are not paraphrases of the sixty-four lines. Attributing the first-step principle to every causal theory would be the straw man. Refusing to articulate any further condition, once the lines are in hand, would be the opposite mistake.

The third objection says that the quantum ledger puts a classical computer back under the table. If the loop is unitary, then the classical record must be ignorance of some deterministic consistent rule. The inclusion that does hold is the dull one. A mixture of process functions can be written with a control register that selects the component and is sent to the sink, and the loop is then unitary because it is block diagonal in that register. The converse is false for one specified case, and the case should be kept small. There is a known quantum board, not a new one, which induces in the computational basis, from a stated source state, a valid classical distribution \(p\) on three binary parties with
\[
L(p) = \frac14 - \frac{3}{16}\sqrt{2} < 0,
\]
while \(L\) is nonnegative on every three-bit process function and hence on every mixture of them. The distribution is tied to that basis. In another basis, a pure process of the same family can induce a mixture. The result excludes one reading of one record: that this computational-basis distribution is ignorance over process functions. It does not exclude classical reality, hidden variables of other kinds, or a basis-independent ontology. It is a pressure test on the hidden-ledger objection, not a second theory of the dossier. A full characterization of which classical processes arise from pure processes remains open. The one strict inclusion used here is that the convex hull of the process functions is a proper subset of the classically recorded processes of this pure-process kind, in the computational basis of the stated example.

One more pass through the room, still invented, after the programmer has printed a record. The engineer says that the printing settles the matter: something computed the answer, so the workshops are a computer, and a possible one, because the printout exists. The archivist should not agree. The printout shows that a finite predicate was evaluated by a search with a first step, on a machine or a person that already had an implementation of its own. The search’s implementation is level B for the search. It is not level B for the cyclic rule. The rule’s level A status is unchanged by the fact that a different machine found its fixed point. If the engineer then says that any device which prints the record implements the workshops, the two-level proposal gives the reply Chalmers’s conditionals already suggest. Implementation is not shared outputs on one run. It is a pattern of physical transitions under variations that were not the variations of this afternoon. The programmer varied candidates in a list. The workshops, if they were a device, would have to vary their own parts. Those are different variations. A metaphysics that counts every sufficient search as the thing searched for will call a timetable a railway.

Leibniz’s question returns here without his answer. A consistent printout is not a reason the workshops exist. The programmer’s success is a reason a certain string was produced by the programmer’s machine. Confusing those reasons is how a modal conclusion gets smuggled in: the string is actual, the string satisfies the rule, therefore a machine of the rule’s shape is possible. The first two clauses can be true while the workshops remain a description. Possibility of the description’s object needs a premise the printout does not contain. The printout is evidence about the search. It is not evidence about the workshops. Nothing else follows.

The workshop, at the end, contains the same two people. The archivist has a list, a ledger, and a proof that the ledger’s canonical extension is unitary under the local tests exactly when the list is a process function. The engineer still does not have a device. A metaphysics of computers that quotes the archivist and dismisses the engineer has treated a permission to compose as a machine. A metaphysics that quotes only the engineer’s habits, and insists that every determinate computation begins at a first laboratory, has treated a feature of ordinary temporal circuits as a law of every admissible composition. The two-level proposal is the refusal of both dismissals. Level A is real, cyclic examples included, and it can be searched when the table is finite and explicit. Level B is real, and it is not settled by A. Between them there is no theorem that carries a consistent table into a possible world. There is only the standing question the engineer asked at the door, which a counting argument does not close.

\appendix
\section{The finite equivalence, in brief}
\label{app}

The appendix fixes the objects named in the text. Proofs and tables are in the unpublished note. Nothing here is a claim to have re-checked a formalization in this draft.

Each party \(k\) has finite nonempty alphabets \(I_k\) and \(O_k\). A local intervention is a tuple of functions \(f_k\colon I_k\to O_k\). A rule \(w\colon O\to I\) is a process function when, for every intervention, the self-map \(o\mapsto (f_k(w_k(o)))_k\) has exactly one fixed point. For the permutation form, identify each \(I_k\) with a cyclic group of the same order. For every rule \(w\), not only for process functions,
\[
U\lvert o\rangle\lvert e\rangle = \lvert w(o)+e\rangle\lvert o\rangle
\]
is a permutation of a finite basis. Given unitaries \(V_k\) from \(I_k\otimes X_k\) to \(O_k\otimes Y_k\) with \(|I_k|\,|X_k|=|O_k|\,|Y_k|\), the closed loop \(M\) has entries
\[
M[(s,y),(e,x)] = \prod_k V_k[(s_k,y_k),(w_k(s)+e_k,x_k)].
\]
If \(w\) is a process function and \(s\neq s'\), some party has \(s_k\neq s'_k\) and \(w_k(s)=w_k(s')\). Otherwise one could define a single intervention with both strings as fixed points. That exclusivity makes \(MM^{\dagger}\) the identity, by factorization of the Gram matrix and a change of variables in each factor where the inputs agree. Square dimensions give unitarity of \(M\). Conversely, if an intervention has \(N\neq 1\) fixed points, the maps \(V_k\lvert i,x\rangle = \lvert f_k(i)+x, i\rangle\) are unitary and the squared norm of one column of \(M\) equals \(N\). So \(M\) is unitary for every such family if and only if \(w\) is a process function. The equivalence is for this \(U\) and these tests.

The separating distribution in Section~\ref{sec:obj} is induced in the computational basis by a known board and a stated source. The value of \(L(p)\) is an exact evaluation of an integer functional. Nonnegativity of \(L\) on the process functions is a finite check. Neither fact is a basis-independent claim.

\bibliographystyle{apalike}
\bibliography{jphil}
\end{document}
